\ExerciseSection

\begin{enumerate}
\def\labelenumi{\arabic{enumi})}
\tightlist
\item
  \begin{enumerate}
  \def\labelenumii{\alph{enumii})}
  \tightlist
  \item
    每个零维空间都是完全正则的。§ 4 习题 26 b) 与 § 5 习题 15 中定义的非正规局部紧空间都是零维的。
  \end{enumerate}
\end{enumerate}

\begin{enumerate}
\def\labelenumi{\alph{enumi})}
\setcounter{enumi}{1}
\item
  拓扑空间 X 称为强零维的，如果对 X 的每个闭子集 A 与 A 的每个邻域 U，存在包含于 U 中的 A 的既开又闭的邻域。每个强零维空间都是正规的。正规空间是强零维的当且仅当它的斯通–切赫紧化（§ 1，习题 7）是全不连通的{[}利用 § 4 习题 17 c){]}。
\item
  设 X 是强零维的豪斯多夫空间，\((\mathbf{U}_{n})\) 是 X 中非空开集的递增序列。试证存在
\end{enumerate}

\[
\left(\mathbf {G} _ {n}\right)
\]

\[
\mathbf {x},
\]

\[
\mathbf {G} _ {n} \subset \mathbf {U} _ {n}
\]

\[
n,
\]

\[
\bigcup_ {n} \mathrm{G} _ {n} = \bigcup_ {n} \mathrm{U} _ {n}
\]

（用归纳法定义 \(G_{n}\)。）进而，若 X 是完备正规的（§ 4，习题 8），则 X 的每个开集都是可数个互不相交的、在 X 中既开又闭的集之并。

\begin{enumerate}
\def\labelenumi{\alph{enumi})}
\setcounter{enumi}{3}
\tightlist
\item
  设 X 是正规空间，且是强零维子空间序列 \((\mathbf{A}_{n})\) 之并。试证 X 是强零维的。{[}设 B、C 是 X 的两个不相交闭子集；用归纳法定义 X 中开集的两个递增序列 \((\mathbf{G}_{n}), (\mathbf{H}_{n})\)，使 \(\overline{G}_{n} \cap \overline{H}_{n} \neq \varnothing\)，\(A_{n} \subset G_{n} \cup H_{n}\)，\(B \cap A_{n} \subset G_{n}\)，\(C \cap A_{n} \subset H_{n}\)。{]}
\end{enumerate}

¶ 2) a) 设 X 是可度量化空间。试证下列性质等价：

\(\alpha)\) 存在度量 \(d\)，与 \(\mathbf{X}\) 的拓扑相容，\(\mathbf{X}\) 关于它是超度量空间（§ 2，习题 4）。

\(\beta)\) X 是强零维的。

γ) 存在 X 的既开又闭子集的局部有限族序列 \((\mathfrak{B}_{n})\)，使 \(B = \bigcup_{n} B_{n}\) 是 X 的拓扑的基。

{[}为证 \(\alpha\) ) 蕴含 \(\beta\) )，注意在超度量空间 X 中，若 F 是闭集，则由所有 \(x \in X\) 中使 \(d(x, F) = \rho > 0\) 者组成的集在 X 中既开又闭。为证 \(\beta\) ) 蕴含 \(\gamma\) )，利用 § 4 第 5 号引理 3 与 § 6 习题 1 c)。最后，为证 \(\gamma\) 蕴含 \(\alpha\) )，化归到 \(\mathfrak{B}_n = (\mathrm{U}_{n,\lambda})_{\lambda \in L}\) 的情形；若 \(f_{n,\lambda}\) 是 \(\mathrm{U}_{n,\lambda}\) 的特征函数，则考察映射 \(x \to f(x) = f_{n,\lambda}(x)\)（X 到乘积群 \(\mathbf{G}^{\mathbf{N} \times \mathbf{L}}\) 中的），其中 \(\mathbf{G} = \mathbf{Z}/2\mathbf{Z}\)（等同于集合 \(\{0, 1\}\)）；对每个 \(n \in \mathbb{N}\)，用 \(\mathbf{G}_n\) 表示 G 中由这样的元素组成的子群：其 \((k, \lambda)\) 坐标对所有 \(\lambda \in L\) 与所有 \(k < n\) 均为零；证明若给 G 赋予以 \((\mathrm{G}_n)\) 为 o 的基本邻域系的群拓扑，则 \(f\) 是 X 到 G 的一个子空间上的同胚；最后注意 G 的拓扑可由一个不变度量定义，G 关于它是超度量空间，从而完成证明。{]}

\begin{enumerate}
\def\labelenumi{\alph{enumi})}
\setcounter{enumi}{1}
\tightlist
\item
  由 a) 推出：每个具有可数基的零维豪斯多夫空间 X 都是强零维的可度量化空间。（利用 § 2 第 8 号命题 13 证明存在由既开又闭的集组成的可数基。）此外，此时 X 同胚于康托尔三分集 K 的一个子空间（参见 § 2 习题 12 (*)。）
\end{enumerate}

¶ 3) a) 试证 R 的每个全不连通子空间都是零维的。

\begin{enumerate}
\def\labelenumi{\alph{enumi})}
\setcounter{enumi}{1}
\tightlist
\item
  设 K 是康托尔三分集；每个 \(x \in K\) 都可唯一地写成形式 \(x = \sum_{k} 2/3^{n_k}\)，其中 \((n_k)_{k \geqslant 1}\) 是严格递增的（有限或无穷的）整数序列 \textgreater0（第 IV 章，§ 8，习题 9）；并定义 \(f(x)\) 为数 \(\sum_{k} (-1)^{n_k}/2^k\)。设 G 是函数 f 在 \(K \times R\) 中的图像。试证可度量化空间 G 是全不连通的但不是零维的。{[}注意 \(\overline{G}\) 与直线 \(\{0\} \times R\) 的交包含一个含点 \((0, 0)\) 的开区间。{]}
\end{enumerate}

\par\noindent\textbf{4) 设 \(\mathbf{X}\) 是可度量化空间。}\par

\begin{enumerate}
\def\labelenumi{\alph{enumi})}
\item
  试证 X 的波莱尔子集的集是 \(\mathfrak{P}(X)\) 中最小的子集 F，它包含 X 的所有闭子集，且 F 中集合的可数并与可数交仍属于 F。（考察 F 的由所有 \(A \in F\) 中使 \(X - A \in F\) 者组成的子集 G。）
\item
  试证 X 的波莱尔子集的集是 \(\mathfrak{P}(X)\) 中最小的子集 F，它包含 X 的所有开子集，且 F 中集合的每个可数交属于 F，F 中互不相交集合的序列的每个并属于 F（同样的方法）。
\item
  每个序数 \(\alpha\) 都可唯一地写成形式 \(\alpha = \omega\beta + n\)，其中 \(n < \omega\)（《集合论》第 III 章，§ 2，习题 15）；\(\alpha\) 称为偶的或奇的，视 \(n\) 是偶数或奇数而定。对每个可数序数 \(\alpha\)，用超限归纳定义子集的集 \(\mathfrak{F}_{\alpha}(\mathbf{X})\)、\(\mathfrak{G}_{\alpha}(\mathbf{X})\)（或简记为 \(\mathfrak{F}_{\alpha}\)、\(\mathfrak{G}_{\alpha}\)）如下：\(\mathfrak{F}_0\)（相应地，\(\mathfrak{G}_0\)）是 X 的所有闭（相应地，开）子集组成的集；若 \(\alpha\) 是偶的且 \(>0\)，则 \(\mathfrak{F}_{\alpha}\)（相应地，\(\mathfrak{G}_{\alpha}\)）是由属于 \(\bigcup_{\xi < \alpha} \mathfrak{F}_{\xi}\)（相应地，\(\bigcup_{\xi < \alpha} \mathfrak{G}_{\xi}\)）的集合的可数交（相应地，可数并）组成的 X 的子集的集；若 \(\alpha\) 是奇的，则 \(\mathfrak{F}_{\alpha}\)（相应地，\(\mathfrak{G}_{\alpha}\)）是由属于 \(\bigcup_{\xi < \alpha} \mathfrak{F}_{\xi}\)（相应地，\(\bigcup_{\xi < \alpha} \mathfrak{G}_{\xi}\)）的集合的可数并（相应地，可数交）组成的 X 的子集的集。试证诸 \(\mathfrak{F}_{\alpha}\)（相应地，\(\mathfrak{G}_{\alpha}\)）的并就是 X 的波莱尔子集的集{[}利用 \(a\){]}。由此推出：若 X 是可数型的，则 X 的波莱尔子集的集的势至多等于连续统的势。
\end{enumerate}

\(d\) ) 关系 \(\mathbf{A}\in\mathfrak{F}_{\alpha}\) 等价于 \(\mathbb{C}\mathbf{A}\in\mathfrak{G}_{\alpha}\)。\(\mathfrak{F}_{\alpha}\)（相应地，\(\mathfrak{G}_{\alpha}\)）中集合的每个有限并（相应地，交）属于 \(\mathfrak{F}_{\alpha}\)（相应地，\(\mathfrak{G}_{\alpha}\)）。我们有 \(\mathfrak{F}_{\alpha}\subset\mathfrak{F}_{\alpha+1}\cap\mathfrak{G}_{\alpha+1}\) 与 \(\mathfrak{G}_{\alpha}\subset\mathfrak{F}_{\alpha+1}\cap\mathfrak{G}_{\alpha+1}\)（用超限归纳）。

\begin{enumerate}
\def\labelenumi{\alph{enumi})}
\setcounter{enumi}{4}
\tightlist
\item
  设 \(f\) 是 \(\mathbf{X}\) 到可度量化空间 \(\mathbf{Y}\) 的连续映射。试证 \(f^{-1}(\mathfrak{F}_{\alpha}(\mathbf{Y})) \subset \mathfrak{F}_{\alpha}(\mathbf{X})\) 且 \(f^{-1}(\mathfrak{G}_{\alpha}(\mathbf{Y})) \subset \mathfrak{G}_{\alpha}(\mathbf{X})\)。
\end{enumerate}

\(f)\) 若 \(\mathbf{X}\) 与 \(\mathbf{Y}\) 是两个可度量化空间，且 \(\mathbf{A} \in \mathfrak{F}_{\alpha}(\mathbf{X}), \mathbf{B} \in \mathfrak{F}_{\alpha}(\mathbf{Y})\)，{[}相应地，\(\mathbf{A} \in \mathfrak{G}_{\alpha}(\mathbf{X}), \mathbf{B} \in \mathfrak{G}_{\alpha}(\mathbf{Y})\){]}，则 \(\mathbf{A} \times \mathbf{B} \in \mathfrak{F}_{\alpha}(\mathbf{X} \times \mathbf{Y})\){[}相应地，\(\mathbf{A} \times \mathbf{B} \in \mathfrak{G}_{\alpha}(\mathbf{X} \times \mathbf{Y})\){]}。

\begin{enumerate}
\def\labelenumi{\alph{enumi})}
\setcounter{enumi}{6}
\item
  设 \((\mathbf{X}_n)\) 是可度量化空间的序列。试证：若 \(\mathbf{A}_n\) 是 \(\mathbf{X}_n\) 中的波莱尔集（对每个 \(n\)），则 \(\prod_{n} \mathbf{A}_n\) 是 \(\prod_{n} \mathbf{X}_n\) 中的波莱尔集。
\item
  设 \(\alpha > 0\) 是可数的偶（相应地，奇）序数，\((\mathbf{A}_n)\) 是 \(\mathbf{X}\) 的可数覆盖，由属于 \(\mathfrak{G}_{\alpha}\)（相应地，\(\mathfrak{F}_{\alpha}\)）的集组成。试证存在分划 \((\mathbf{B}_n)\)（\(\mathbf{X}\) 的），使 \(\mathbf{B}_n \subset \mathbf{A}_n\) 且 \(\mathbf{B}_n \in \mathfrak{F}_{\alpha} \cap \mathfrak{G}_{\alpha}\)（对每个 \(n\)）{[}利用 \(d\){]}。
\end{enumerate}

¶ 5) a) 设 J 是波兰空间 \(\mathbf{N}^{\mathbf{N}}\)（N 赋有离散拓扑），X 是可数型的可度量化空间。试证：对每个可数序数 \(\alpha\)，存在子集 \(\mathbf{G}_{\alpha}\)（\(\mathbf{J} \times \mathbf{X}\) 的），使得 (i) \(\mathbf{G}_{\alpha} \in \mathfrak{G}_{\alpha}\)（\(\mathbf{J} \times \mathbf{X}\)）（习题 4），且 (ii) 对每个子集 \(\mathbf{U} \in \mathfrak{G}_{\alpha}(\mathbf{X})\)，存在 \(z \in \mathbf{J}\) 使 \(\mathbf{G}_{\alpha}(z) = \mathbf{U}\)。{[}可用如下的超限归纳进行：空间 J 与 \(\mathbf{J}^{\mathbf{N}}\) 同胚，设 \(f\) 是 J 到 \(\mathbf{J}^{\mathbf{N}}\) 上的同胚，对每个整数 \(n\)，令 \(f_{n} = \mathrm{pr}_{n} \circ f\)。若 \((\mathbf{A}_{n})\) 是 X 的包含空集的可数基，取 \(\mathbf{G}_{0}\) 使 \(\mathbf{G}_{0}(z) = \bigcup_{n} \mathbf{A}_{f_{n}(z)}\)（对每个 \(z \in \mathbf{J}\)）。对每个可数序数 \(\alpha\)，取 \(\mathbf{G}_{\alpha+1}\) 使 \(\mathbf{G}_{\alpha+1}(z) = \bigcap_{n} \mathbf{G}_{\alpha}(f_{n}(z))\)（若 \(\alpha\) 为偶），而 \(\mathbf{G}_{\alpha+1}(z) = \bigcup_{n} \mathbf{G}_{\alpha}(f_{n}(z))\)（若 \(\alpha\) 为奇）。最后，若 \(\alpha\) 无前驱且 \((\lambda_{n})\) 是满足 \(\alpha = \sup_{n} \lambda_{n}\) 的递增序数序列，取 \(\mathbf{G}_{\alpha}\) 使 \(\mathbf{G}_{\alpha}(z) = \bigcup_{n} \mathbf{G}_{\lambda_{n}}(f_{n}(z))\)。利用 \(f_{n}\) 在 J 上连续这一事实。{]}

\begin{enumerate}
\def\labelenumi{\alph{enumi})}
\setcounter{enumi}{1}
\tightlist
\item
  特别地取 \(\mathbf{X} = \mathbf{J}\)。试证集 \(\mathrm{pr}_1(\mathbf{G}_\alpha \cap \Delta)\)（其中 \(\Delta\) 是 \(\mathbf{J} \times \mathbf{J}\) 的对角线）属于 \(\mathfrak{G}_{\alpha}(\mathbf{J})\) 但不属于 \(\mathfrak{F}_{\alpha}(\mathbf{J})\){[}用反证法，利用 \(a\){]}。
\end{enumerate}

\begin{enumerate}
\def\labelenumi{\arabic{enumi})}
\setcounter{enumi}{5}
\tightlist
\item
  \begin{enumerate}
  \def\labelenumii{\alph{enumii})}
  \tightlist
  \item
    设 J 是波兰空间 \(N^{N}\)。试证：若 X 是任一苏斯林空间，则存在 J 到 X 上的连续满射{[}化归到 X 为波兰空间的情形，并考察用一个所有 \(C_{n}\) 皆为无穷的筛 (\(C_{n}\)) 对 X 所作的筛分{]}。
  \end{enumerate}
\end{enumerate}

\begin{enumerate}
\def\labelenumi{\alph{enumi})}
\setcounter{enumi}{1}
\item
  设 X 是可度量化空间。试证：若 A 是 X 的任一苏斯林子空间，则存在 \(J \times X\) 的闭子集 F，使 \(\mathbf{A} = \operatorname{pr}_{2}(\mathbf{F})\){[}利用 a){]}。
\item
  设 X 是可数型的可度量化空间；令 \(L = J \times X\)，并设 F 是 \(J \times L\) 的闭子集，使得对每个闭子集
\end{enumerate}

M，存在 \(z \in J\) 使 \(F(z) = M\){[}习题 5 a){]}。试证：对 X 的每个苏斯林子空间 S，存在 \(z \in J\) 使 \(\operatorname{pr}_{2}(\mathbf{F}(z)) = \mathbf{S}\)。由此推出：若 X = J，则由所有 \(z \in J\) 中使 \(z \in \operatorname{pr}_{2}(\mathbf{F}(z))\) 者组成的集 T 是苏斯林集，但 J --- T 不是苏斯林集{[}推理与习题 5 b) 相同{]} (*)。

\begin{enumerate}
\def\labelenumi{\arabic{enumi})}
\setcounter{enumi}{6}
\tightlist
\item
  \begin{enumerate}
  \def\labelenumii{\alph{enumii})}
  \tightlist
  \item
    试证每个零维波兰空间都同胚于乘积 \(J = N^{N}\) 的一个闭子空间（考察用既开又闭的集对 X 所作的一个严格筛分）。
  \end{enumerate}
\end{enumerate}

\begin{enumerate}
\def\labelenumi{\alph{enumi})}
\setcounter{enumi}{1}
\item
  设 X 是零维波兰空间，Y 是 X 的稠密子空间，它没有内点且是 X 中开集的可数交。试证 Y 同胚于 J。（注意在可度量化的零维空间中，开而不闭的集是可数无穷个两两不相交的既开又闭的集之并；利用第 1 号定理 1。）由此推出：R 的每个没有内点且为开集的可数交的稠密子空间都同胚于 J（注意这样的集包含在 R 的某个可数稠密子集 D 的补集中，且 CD 是零维的）。
\item
  若 X 是任一不可数的零维波兰空间，则存在把 X 分为一个可数集与一个同胚于 J 的子空间的分划{[}利用 b) 与 § 5 习题 11{]}。
\end{enumerate}

\begin{enumerate}
\def\labelenumi{\arabic{enumi})}
\setcounter{enumi}{7}
\tightlist
\item
  设 X 是苏斯林空间，f 是 X 到豪斯多夫空间 Y 的连续映射。试证：若 \(f(\mathbf{X})\) 不可数，则存在 X 的同胚于康托尔三分集 K（第 IV 章，§ 2，第 5 号）的子空间 A，使 \(f|A\) 是单射。{[}化归到 X 是完备度量空间的情形：证明存在筛 \(\mathbf{C} = (\mathbf{C}_{n}, p_{n})\)，使对每个 n，\(C_{n}\) 有 \(2^{n}\) 个元素，且对每个 n 存在映射 \(\varphi_{n}\)（从 \(C_{n}\) 到 X 中直径 \(\leqslant 2^{-n}\) 的非空闭子集全体之集），使得 (i) \(\varphi_{n+1}(c) \subset \varphi_{n}(p_{n}(c))\)（对所有 \(c \in C_{n+1}\)）；(ii) 每当 c 与 \(c'\) 是 \(C_{n}\) 的不同元素时，\(f(\varphi_{n}(c))\) 与 \(f(\varphi_{n}(c'))\) 不相交。利用 § 5 习题 11。{]}特别地，每个不可数的苏斯林空间都含有同胚于 K 的子空间，因而具有连续统的势。
\end{enumerate}

¶ 9) a) 设 X 是波兰空间，R 是 X 上的等价关系。试证：在下列两个假设之一成立时，存在 X 的一个子集，它是开集的可数交，且与每个等价类恰交于一点：

\(\alpha)\) R 是闭的；

\(\beta)\) \(\mathbf{X}\) 局部紧且 \(\mathbf{R}\) 的图像在 \(\mathbf{X} \times \mathbf{X}\) 中闭。

{[}对 \(\alpha\) )，沿用第 8 号定理 4 的证明方法，并利用如下注记：若 \((\mathbf{G}_n)\) 是这样的集合序列，其中每个都是开集的可数交，且对每个 \(n\) 存在邻域 \(\mathbf{V}_n\)（\(\mathbf{G}_n\) 的）使族 \((\mathbf{V}_n)\) 局部有限，则 \(\bigcup_{n} \mathbf{G}_n\) 是开集的可数交。对 \(\beta\) ) 用同样的方法；

注意紧集的饱和化现在是闭的（参见第 I 章 § 10 习题 16）。{]}

\begin{enumerate}
\def\labelenumi{\alph{enumi})}
\setcounter{enumi}{1}
\item
  设 X 是波兰空间，R 是 X 上的既开又闭的等价关系，\(\varphi: \mathbf{X} \to \mathbf{X} / \mathbf{R}\) 是典范映射。试证存在 X 的子集 A，它是 X 中开集的可数交，使得 \(\varphi\) 在 A 上的限制是 A 到 \(\varphi(A)\) 上的同胚，且 \(\varphi(A)\) 在 X/R 中稠密并是开集的可数交。{[}用第 8 号定理 4 证明中的记号，证明可以假设对每个 \(n\)，诸 \(\varphi_n(c)\) 已经如此定义，使得 \(\varphi\) 作用于诸 \(h_n(c)\) 的内部之并（当 \(c\) 取遍 \(C_n\) 时）所得的像在 X/R 中稠密；利用 \(a\) ) 与第 1 号定理 1。{]}
\item
  试证当 X 局部紧且具有可数基，而 R 是 X 上的闭等价关系且其所有等价类都紧时，b) 的结论仍成立。{[}利用 § 5 习题 26 化归到情形 b)。{]}
\end{enumerate}

¶ 10) a) 设 X 是可度量化空间，S 是 X 的苏斯林子空间。试证 S 在 X 中是几乎开集（§ 5，习题 6）。{[}设 P 是波兰空间，g 是 P 到 S 上的连续映射，(Cn, pn, φn) 是 P 的一个筛分，f 是 L(C) 到 P 上的相应连续映射（用第 5 号的记号）；令 h = g∘f，对每个 c ∈ Cn，用 qn(c) 表示由满足 cn = c 的序列 (ck) 组成的 L(C) 的子空间。对每个 c ∈ Cn，令 Fn(c) = h(qn(c))，并令 Un(c) = g(φn(c)) ∃ Fn(c)；用 Wn(c) 表示 D(Un(c))（§ 5，习题 3）与 X 中一个贫集之并，使 Fn(c) ⊂ Wn(c) ⊂ U̅n(c)，从而 Wn(c) 在 X 中是几乎开集。证明

\[
\mathrm{V} _ {n} (c) = \mathrm{W} _ {n} (c) \cap \mathbb {C} \left(\bigcup_ {p _ {n} \left(c ^ {\prime}\right) = c} \mathrm{W} _ {n + 1} \left(c ^ {\prime}\right)\right)
\]

在 X 中是贫集，注意 \(\mathbf{F}_n(c) = \bigcup_{p_n(c') = c}\mathbf{F}_{n + 1}(c')\)；再证明集 \(\mathrm{W}_{n}(c)-\mathrm{F}_{n}(c)\) 包含在诸集 \(\mathrm{V}_{m}(d)\) 之并中，其中 \(m\geqslant n\)，且对每个 m，d 取遍 \(C_{m}\) 中满足 \(p_{nm}(d)=c\) 的元素。由此证明 \(\mathrm{F}_{n}(c)\) 是几乎开集。{]}

\begin{enumerate}
\def\labelenumi{\alph{enumi})}
\setcounter{enumi}{1}
\tightlist
\item
  给出连续映射 \(f\)（从 {[}0, 1{]} 到自身）与几乎开集 \(\mathbf{B} \subset \mathbf{I}\) 的一个例子，使 \(f(\mathbf{B})\) 不是几乎开集。{[}利用第 IV 章 § 8 习题 16 b)，并证明可取 \(\mathbf{B}\) 为康托尔集 \(\mathbf{K}\) 的这样的子集，使 \(f(\mathbf{B}) = \mathbf{Z}\) 具有 § 5 习题 18 d) 中描述的性质。{]}
\end{enumerate}

¶ 11) 设 \(\mathbf{X} = \mathbf{R}^n\)，\(\mathbf{M}\) 是 \(\mathbf{X}\) 的闭子集。点 \(x \in \mathbf{M}\) 称为线性可达的，如果存在点 \(y \in \mathbf{X} - \mathbf{M}\)，使以 \(x, y\) 为端点的开线段包含于 \(\mathbf{X} - \mathbf{M}\)。试证集 \(\mathbf{L} \subset \mathbf{M}\)（由 \(\mathbf{M}\) 的线性可达点组成）是苏斯林集。{[}设 \(d(x, y)\) 是 \(\mathbf{X}\) 中的欧氏度量，对每个 \(z \in \mathbf{X}\)，令 \(f_z(x, y) = d(x, z) + d(z, y) - d(x, y)\)。注意所有 \((x, y, z) \in \mathbf{X} \times \mathbf{X} \times \mathbf{X}\) 中使 \(x \in \mathbf{M}, y \in \mathbf{X} - \mathbf{M}, z \in \mathbf{M}, z \neq x\) 且 \(f_z(x, y) = 0\) 者组成的集是紧集的可数并，因而它在前两个因子的乘积上的投影也是紧集的可数并。{]}

\begin{enumerate}
\def\labelenumi{\arabic{enumi})}
\setcounter{enumi}{11}
\tightlist
\item
  设 X 是可度量化空间，f 是 X 的紧子集的集 \(\mathcal{R}(\mathbf{X})\) 到 \(\overline{R}\) 中的映射，满足条件 \((\mathrm{CA}_{\mathrm{I}})\)（对 X 的两个紧子集）与 \((\mathrm{CA}_{\mathrm{III}})\)。对 X 的每个子集 A，用 \(f_{*}(\mathrm{A})\) 表示对所有紧子集 K ⊂ A 的数 \(f(\mathrm{K})\) 的上确界，用 \(f^{*}(\mathrm{A})\) 表示对所有包含 A 的开集 U 的 \(f_{*}(\mathrm{U})\) 的下确界。若 \(f_{*}(\mathrm{A}) = f_{*}(\mathrm{A})\)，则称 A 关于 f 是容许的。
\end{enumerate}

\begin{enumerate}
\def\labelenumi{\alph{enumi})}
\item
  试证：对 X 的每个紧子集 L 与每个 \(\varepsilon > 0\)，存在开集 \(U \supset L\)，使对满足 \(L \subset K \subset U\) 的 X 的每个紧子集 K，有 \(f(K) \leqslant f(L) + \varepsilon\)。（用反证法。）由此推出 X 的每个紧子集与每个开子集关于 f 都是容许的。
\item
  假设 \(f\) 满足下列条件：
\end{enumerate}

(AL) 对 X 的每对紧子集 K, K'，有

\[
f (\mathbf {K} \cup \mathbf {K} ^ {\prime}) + f (\mathbf {K} \cap \mathbf {K} ^ {\prime}) \leqslant f (\mathbf {K}) + f (\mathbf {K} ^ {\prime}).
\]

试证：若 \((\mathbf{K}_i)\) 是具有下列性质的紧集的有限族：

\[
f \left(\bigcup_ {i} \mathrm{K} _ {i}\right) <   + \infty ,
\]

且 \((\mathbf{K}_i')\) 是以同一指标集为指标的另一紧集有限族，使 \(\mathbf{K}_i' \subset \mathbf{K}_i\) 且 \(f(\mathbf{K}_i') > -\infty\)（对每个 \(i\)），则

\[
f \left(\bigcup_ {i} \mathrm{K} _ {i}\right) - f \left(\bigcup_ {i} \mathrm{K} _ {i} ^ {\prime}\right) \leqslant \sum_ {i} \left(f \left(\mathrm{K} _ {i}\right) - f \left(\mathrm{K} _ {i} ^ {\prime}\right)\right).
\]

设 \((\mathbf{A}_i), (\mathbf{B}_i)\) 是 \(\mathbf{X}\) 的子集的两个有限族，具有同一指标集，使 \(\mathbf{B}_i \subset \mathbf{A}_i\)（对每个 \(i\)），\(f^* \left( \bigcup_i \mathbf{A}_i \right) < +\infty\) 且 \(f^*(\mathbf{B}_i) > -\infty\)（对每个 \(i\)）。试证：若 \(f\) 满足 (AL)，则

\[
f ^ {*} \left(\bigcup_ {i} \mathrm{A} _ {i}\right) - f ^ {*} \left(\bigcup_ {i} \mathrm{B} _ {i}\right) \leqslant \sum_ {i} \left(f ^ {*} (\mathrm{A} _ {i}) - f ^ {*} (\mathrm{B} _ {i})\right).
\]

{[}化归到指标集含两个元素的情形。先考虑 \(A_{i}\) 与 \(B_{i}\) 都是开集的情形，并利用下面的引理：若 \(U_{1}, U_{2}\) 是两个开集，K 是包含于 \(U_{1} \cup U_{2}\) 中的紧集，则存在紧集 \(K_{i} \subset U_{i}\)（i = 1, 2），使 \(K \subset K_{1} \cup K_{2}\)。{]} c) 假设 f 满足条件 (AL)。试证：对 X 的子集的每个递增序列 \((\mathbf{A}_{n})\)（满足 \(f^{*}(\mathbf{A}_{n}) > -\infty\)），有 \(f^{*}\left(\bigcup_{n} \mathbf{A}_{n}\right) = \sup_{n} f^{*}(\mathbf{A}_{n})\){[}利用 b){]}。由此推出：若 f 不取值 \(-\infty\)（在 \(\Re(\mathbf{X})\) 上），则 \(f^{*}\) 是 X 上的一个容度。

\begin{enumerate}
\def\labelenumi{\alph{enumi})}
\setcounter{enumi}{3}
\tightlist
\item
  假设 \(f\) 满足条件 (AL)。试证：任意容许集序列 \((\mathbf{A}_n)\)（使 \(f^*(\mathbf{A}_n) > -\infty\)，对每个 \(n\)）之并是容许的。{[}利用 b) 与 c)。{]}
\end{enumerate}

* 13) 设 K 是 \(\mathbf{R}^2\) 的紧子集。对每个实数 \(y\)，用 \(\delta_1(\mathbf{K}, y)\){[}相应地，\(\delta_2(\mathbf{K}, y)\){]}表示 \(\mathbf{K} \cap ([0, +\infty [\times \{y\})]\){[}相应地，\(\mathbf{K} \cap (]-\infty, 0] \times \{y\})\){]}的直径。试证函数 \(y \to \delta_1(\mathbf{K}, y)\) 与 \(y \to \delta_2(\mathbf{K}, y)\) 都是上半连续的。设 \(\varphi\) 是 \([0, +\infty]\) 到自身的递增连续映射，满足 \(\varphi(0) = 1\) 与 \(\varphi(+\infty) = 2\)，并令

\[
\psi (\mathbf {K}, y) = \varphi (\delta_ {1} (\mathbf {K}, y) \delta_ {2} (\mathbf {K}, y)) \quad \text { and } \quad f (\mathbf {K}) = \int_ {\mathrm{pr} _ {2} (\mathbf {K})} \psi (\mathbf {K}, y) d y.
\]

试证 \(f\) 满足条件 \((\mathbf{CA}_{\mathbf{I}})\) 与 \((\mathbf{CA}_{\mathbf{III}})\)，且有 \(f(\mathbf{K} \cup \mathbf{K}') \leqslant f(\mathbf{K}) + f(\mathbf{K}')\)，对紧子集 \(\mathbf{K}, \mathbf{K}'\)（\(\mathbf{R}^2\) 中的）的每一对成立。但若 \(\mathbf{A}\) 是由 \(x \geqslant 0\) 与 \(0 \leqslant y \leqslant 1\) 定义的闭集，试证 \(f_*(\mathbf{A}) = 1\) 且 \(f^*(\mathbf{A}) = 2\)。

\begin{enumerate}
\def\labelenumi{\arabic{enumi})}
\setcounter{enumi}{13}
\tightlist
\item
  设 \(\mathbf{X}\) 是可度量化空间，设 \(f\) 是 \(\mathbf{X}\) 上的容度，满足 \(f(\mathbf{A} \cup \mathbf{B}) \leqslant f(\mathbf{A}) + f(\mathbf{B})\)。
\end{enumerate}

\begin{enumerate}
\def\labelenumi{\alph{enumi})}
\item
  用习题 12 的记号，试证 \(f^{*}(\mathbf{A} \cup \mathbf{B}) \leqslant f^{*}(\mathbf{A}) + f^{*}(\mathbf{B})\) 与 \(f_{*}(\mathbf{A} \cup \mathbf{B}) \leqslant f_{*}(\mathbf{A}) + f^{*}(\mathbf{B})\)，对任意两个子集 \(\mathbf{A}, \mathbf{B}\)（均属于 \(\mathbf{X}\)）成立{[}若 \(\mathbf{K}\) 是包含于 \(\mathbf{A} \cup \mathbf{B}\) 中的紧集，\(\mathbf{U}\) 是包含 \(\mathbf{B}\) 的开集，把 \(\mathbf{K}\) 写成 \(\mathbf{K} = (\mathbf{K} \cap \mathbf{U}) \cup (\mathbf{K} \cap \mathbf{C}\mathbf{U})\) 的形式。{]}
\item
  设 K 是 X 的具有连续统的势的紧子集，(A, B) 是把 K 分成两个集的一个分划，使 A 或 B 的每个紧子集都是可数的{[}参见 § 5 习题 18 d){]}。试证：若 \(f(\{x\}) = 0\)（对所有 \(x \in \mathbf{X}\)），且 \(f(\mathbf{K}) > 0\)，则 \(\mathbf{A}\) 与 \(\mathbf{B}\) 关于 \(f\) 都不是可容的。
\end{enumerate}

* 15) 设 \(\mu\) 是 \(\mathbf{R}\) 上的勒贝格测度；定义容度 \(f\)（\(\mathbf{R}^2\) 上的，满足习题 12 中条件 (AL)）：令 \(f(\mathbf{A}) = \mu^{*}(\mathrm{pr}_{1}\mathbf{A})\)（命题 15）。

\begin{enumerate}
\def\labelenumi{\alph{enumi})}
\item
  设 A 是 \(R^{2}\) 的不可容有界子集（习题 14），\(B_{0}\) 是圆周 \(\|x\| = r\)，使 A 包含于开圆盘 \(\|x\| < r\) 中，并设 \(B_{1}\) 是圆周 \(\|x\| = r'\)，其半径 \(r' > r\)。试证 \(A \cup B_{0}\) 与 \(A \cup B_{1}\) 都是可容的，尽管其交 A 不是。
\item
  对每个 \(n\)，用 \(\mathbf{C}_n\) 表示由所有满足下式的 \(x \in \mathbb{R}^2\) 组成的集：
\end{enumerate}

\[
r <   | | \boldsymbol {x} | | <   r + \frac {\mathrm{I}}{n}.
\]

试证 \(\mathbf{A} \cup \mathbf{C}_n\) 的每一个都是可容的，尽管它们的交不是，且有 \(\inf_n f(\mathbf{C}_n) \neq f\left(\bigcap_n \mathbf{C}_n\right)\)。 *

¶ 16) 设 X, X' 是两个可度量化空间。映射 \(f: \mathbf{X} \to \mathbf{X}'\) 称为波莱尔映射，如果对 X' 的每个闭子集 F'，\(f^{-1}(\mathbf{F}')\) 是 X 中的波莱尔集。对每个可数的偶（相应地，奇）序数 \(\alpha\)，称 \(f\) 是类 \(\alpha\) 的，如果对 X' 的每个闭子集 F'，\(f^{-1}(\mathbf{F}')\) 属于 \(\mathfrak{S}_{\alpha}(\mathbf{X})\){[}相应地 \(\mathfrak{G}_{\alpha}(\mathbf{X})\){]}（习题 4）；于是，对 X' 的每个开集 G'，\(f^{-1}(\mathbf{G}')\) 属于 \(\mathfrak{G}_{\alpha}(\mathbf{X})\){[}相应地 \(\mathfrak{F}_{\alpha}(\mathbf{X})\){]}。类 0 的波莱尔映射恰好就是连续映射。

\begin{enumerate}
\def\labelenumi{\alph{enumi})}
\item
  特征函数 \(\varphi_{\mathbf{A}}\)（\(\mathbf{A}\) 是 \(\mathbf{X}\) 的子集）是类 \(\alpha\) 的，当且仅当 \(\mathbf{A} \in \mathfrak{F}_{\alpha} \cap \mathfrak{G}_{\alpha}\)。
\item
  设 \(f: \mathbf{X} \to \mathbf{X}'\) 是类 \(\alpha\) 的映射。试证：对每个子集 \(\mathbf{B}' \in \mathfrak{F}_{\beta}(\mathbf{X}')\){[}相应地 \(\mathbf{B}' \in \mathfrak{G}_{\beta}(\mathbf{X}')\){]}，\(\overline{f}(\mathbf{B}')\) 属于 \(\mathfrak{F}_{\alpha + \beta}(\mathbf{X})\){[}相应地 \(\mathfrak{G}_{\alpha + \beta}(\mathbf{X})\){]}。（对 \(\beta\) 用超限归纳法。）
\item
  设 \(f: \mathbf{X} \to \mathbf{X}'\) 是类 \(\alpha\) 的映射，\(g: \mathbf{X}' \to \mathbf{X}''\) 是类 \(\beta\) 的映射，其中 \(\mathbf{X}''\) 是第三个可度量化空间。试证 \(g \circ f\) 是类 \(\alpha + \beta\) 的波莱尔映射。
\item
  对每个可数的偶（相应地，奇）序数 \(\alpha\)，设 \((\mathbf{A}_n)\) 是 \(\mathbf{X}\) 的子集的序列，这些子集属于 \(\mathfrak{G}_{\alpha}\)（相应地，\(\mathfrak{F}_{\alpha}\)）且覆盖 \(\mathbf{X}\)。若映射 \(f: \mathbf{X} \to \mathbf{X}'\) 使 \(f|_{\mathbf{A}_n}\) 是类 \(\alpha\) 的（对每个 \(n\)），则 \(f\) 是类 \(\alpha\) 的。对 \(\mathbf{X}\) 的子集的属于 \(\mathfrak{F}_{\alpha}\)（相应地，\(\mathfrak{G}_{\alpha}\)）且覆盖 \(\mathbf{X}\) 的有限序列，给出类似的结果。
\end{enumerate}

\begin{enumerate}
\def\labelenumi{\arabic{enumi})}
\setcounter{enumi}{16}
\tightlist
\item
  \begin{enumerate}
  \def\labelenumii{\alph{enumii})}
  \tightlist
  \item
    设 \(\mathbf{X}, \mathbf{X}'\) 是两个可度量化空间，其中 \(\mathbf{X}'\) 是可数型的，且设 \((\mathbf{U}_n')\) 是 \(\mathbf{X}'\) 的拓扑的可数基。设 \(\alpha\) 是可数的偶（相应地，奇）序数。若映射 \(f: \mathbf{X} \to \mathbf{X}'\) 使 \(f^{-1}(\mathbf{U}_n')\) 属于 \(\mathfrak{G}_{\alpha}(\mathbf{X})\){[}相应地 \(\mathfrak{F}_{\alpha}(\mathbf{X})\){]}（对每个 \(n\)），则 \(f\) 是类 \(\alpha\) 的波莱尔映射。
  \end{enumerate}
\end{enumerate}

\begin{enumerate}
\def\labelenumi{\alph{enumi})}
\setcounter{enumi}{1}
\item
  设 \((\mathbf{X}_{n}^{\prime})\) 是可数型的可度量化空间的序列。映射 \(f = (f_{n}) : \mathbf{X} \to \prod_{n} \mathbf{X}_{n}^{\prime}\) 是类 \(\alpha\) 的，当且仅当每个 \(f_{n}\) 都是类 \(\alpha\) 的{[}利用 a){]}。由此推出：若 X 是可数型的，则 X 上类 \(\alpha\) 的有限实值函数组成一个环。
\item
  设 X, X' 是两个可数型的可度量化空间，且 \(\alpha\) 是可数的偶（相应地，奇）序数。若映射 \(f: X \to X'\) 是类 \(\alpha\) 的，试证其图像属于 \(\mathfrak{F}_{\alpha}(X \times X')\){[}相应地 \(\mathfrak{G}_{\alpha}(X \times X')\){]}{[}利用 b) 与习题 16 b){]}。
\end{enumerate}

\begin{enumerate}
\def\labelenumi{\arabic{enumi})}
\setcounter{enumi}{17}
\tightlist
\item
  设 X 是苏斯林空间，\(X'\) 是可数型的可度量化空间，f 是 X 到 \(X'\) 中的映射。
\end{enumerate}

\begin{enumerate}
\def\labelenumi{\alph{enumi})}
\item
  试证：若 f 的图像 G 是 \(X \times X'\) 中的苏斯林集，则 f 是波莱尔函数，因而 G 是波莱尔集{[}习题 17 c){]}（利用定理 2 的推论）。
\item
  试证：若 \(f\) 是单射波莱尔映射且 \(\mathbf{X}'\) 是苏斯林空间，则任何苏斯林集在 \(f\) 下的像（该集含于 \(\mathbf{X}\)）是 \(\mathbf{X}'\) 中的苏斯林集（利用习题 17c)）。若 \(f\) 是双射的，则逆映射也是波莱尔映射。
\end{enumerate}

\begin{enumerate}
\def\labelenumi{\arabic{enumi})}
\setcounter{enumi}{18}
\tightlist
\item
  设 X, X' 是两个可度量化空间，f 是 X 到 X' 中的几乎开映射（§ 5，习题 24）。试证：若 B' 是 X' 中的任一波莱尔集，则 \(f^{-1}(B')\) 在 X 中是几乎开的（参见 § 5，习题 6）。由此推出：若 g 是 X' 到可度量化空间 X'\,中的波莱尔映射，则 g∘f 是几乎开的。给出连续映射 f 与几乎开映射 g 的一个例子，使 g∘f 不是几乎开的{[}参见习题 10 b){]}。
\end{enumerate}

